\let\negmedspace\undefined
\let\negthickspace\undefined
\documentclass[journal]{IEEEtran}
\usepackage[a5paper, margin=10mm, onecolumn]{geometry}
%\usepackage{lmodern} % Ensure lmodern is loaded for pdflatex
\usepackage{tfrupee} % Include tfrupee package

\setlength{\headheight}{1cm} % Set the height of the header box
\setlength{\headsep}{0mm}     % Set the distance between the header box and the top of the text

\usepackage{gvv-book}
\usepackage{gvv}
\usepackage{cite}
\usepackage{amsmath,amssymb,amsfonts,amsthm}
\usepackage{algorithmic}
\usepackage{graphicx}
\usepackage{textcomp}
\usepackage{xcolor}
\usepackage{txfonts}
\usepackage{listings}
\usepackage{enumitem}
\usepackage{mathtools}
\usepackage{gensymb}
\usepackage{comment}
\usepackage[breaklinks=true]{hyperref}
\usepackage{tkz-euclide} 
\usepackage{listings}
% \usepackage{gvv}                                        
\def\inputGnumericTable{}                                 
\usepackage[latin1]{inputenc}                                
\usepackage{color}                                            
\usepackage{array}                                            
\usepackage{longtable}                                       
\usepackage{calc}                                             
\usepackage{multirow}                                         
\usepackage{hhline}                                           
\usepackage{ifthen}                                           
\usepackage{lscape}
\usepackage{multicol}
\begin{document}

\bibliographystyle{IEEEtran}
\vspace{3cm}

\title{2.0.12}
\author{EE25BTECH11012-BEERAM MADHURI}
% \maketitle
% \newpage
% \bigskip
{\let\newpage\relax\maketitle}

\renewcommand{\thefigure}{\theenumi}
\renewcommand{\thetable}{\theenumi}
\setlength{\intextsep}{10pt} % Space between text and floats


\numberwithin{equation}{enumi}
\numberwithin{figure}{enumi}
\renewcommand{\thetable}{\theenumi}


\textbf{Question}:\\
A product is made by mixing three raw materials 1, 2, 3 in varying proportions, where material 1 must account for not more than 50\% of the total. If $x, y$ and $z$ are the amounts of raw materials 1, 2, 3, respectively, this constraint can be modeled as \hfill(2007 PI)

\begin{enumerate}[label=\alph*)]
    \item $x \le 0.5$
    \item $x \le 0.5(x + y + z)$
    \item $0.5x \le x + y + z$
    \item $x \ge 0.5(y + z)$
\end{enumerate}\\

\textbf{Solution:}\\
Given,\\
Raw Materials: 1,2,3\\
Amount: x,y,z\\
let,\\
Total amount of Product: T\\
\begin{align}
    T = x+y+z
\end{align}
given constraint: Material one (x) must not account for more than 50\% of the total.\\
\begin{align} 
    x \le 0.5(T)\\
    x \le 0.5(x+y+z)\\
    x\le y+z
\end{align}

Correct Option: b

\begin{figure}
\centering
\includegraphics[width=1\columnwidth]{Figures/plot.png}
\caption{Boundary Condition}
\label{fig:placeholder}
\end{figure}
\end{document}
